\section{Track II: insertion status freeze}
\label{sec:O-ins-freeze}

\paragraph{Retained objects.}
Seed $F_t=\langle t_D,h\otimes h\otimes\phi\rangle$ (XVI-I; no extra factorial) and the
regular, seed-preserving, antifield-linear class
$I_t=F_t+A_0+\kappa A_1$ with $A_0=\int(h^*U+\phi^*V+c^*W)$.

\paragraph{Frozen negative / conditional result.}
Under XVI-N massless-obstruction hypotheses, on the constructed massless special-family
witness: no such $I_t$ satisfies $Q_{\mathrm{inv}}I_t=O(\kappa^2)$. Route~A is sufficient
and discharged; Route~B is not claimed discharged. Certified small-mass interval
$0\le m<1.3715\times10^{-4}$ at frozen geometry; arbitrary prescribed $m>0$ unresolved.

\paragraph{Hard rule.}
No quantum completion is claimed whose classical leading term contradicts this obstruction.
Extra antifields or nonminimal fields do not automatically remove the order-$\kappa^0$
field-only contradiction ($P_0 Q_0 A_0=0$ on the free shell).

\paragraph{O.9 separation.}
Existence of a renormalized product of the bare $F_t$ is not closure of $F_t$ in the
BV/BRST sense. N-Ins.5 remains structural only after invariant completion.
Success on the bulk quantum sector (verdict~1) does not complete $F_t$ (verdict~3).
